% Methodology section for the SecureOS report: file system.
% Drop-in for the report body. If the report already has a \section{Methodology},
% delete the \section line below and keep the \subsection.

\section{Methodology}
\subsection{File System}

The file system was built bottom up, as four layers with a single direction of
dependency: a block device abstraction, the FAT32 on-disk structures, the file
and path operations built on those structures, and finally the system call
layer that exposes them to user programs. Each layer was written and verified
before the one above it was started, so that a failure at any point in
development could be attributed to the layer being worked on rather than to
the ones underneath. Table~\ref{tab:fs-layers} lists the layers and the source
files that implement them.

\begin{table}[H]
    \centering
    \begin{tabular}{|l|l|l|}
        \hline
        \textbf{Layer} & \textbf{Responsibility} & \textbf{Source} \\
        \hline
        Block device & Sector-granular read and write & \texttt{block\_device.c} \\
        Device driver & ATA/IDE transfers in PIO mode & \texttt{ide\_pio.c} \\
        Volume & Boot sector parsing, mount validation & \texttt{fat.c} \\
        Allocation & FAT entries, cluster chains, FSInfo & \texttt{fat32\_cluster.c} \\
        Directory & 8.3 entries, lookup, create, delete & \texttt{fat32\_dir.c} \\
        File and path & Open, read, write, truncate, seek & \texttt{fat32\_fs.c} \\
        System calls & Per-process descriptors, validation & \texttt{fs\_syscalls.c} \\
        \hline
    \end{tabular}
    \caption{File system layers and their implementation files.}
    \label{tab:fs-layers}
\end{table}

\subsubsection{Choice of On-Disk Format}

FAT32 was selected over a purpose-designed format for three reasons. First, it
is fully documented and stable, so the on-disk layout could be treated as a
fixed specification to implement against rather than a moving design decision.
Second, standard host tooling can produce and inspect a FAT32 volume, which
means the disk images used during development are created by an independent
implementation (\texttt{mkfs.fat}) and populated by another (\texttt{mcopy});
any disagreement between our reader and the on-disk bytes is therefore a bug in
our reader, not an ambiguity shared by both sides. Third, FAT32 has no journal
and no extent trees, so a correct implementation is small enough to be read and
audited in full, which is consistent with the project's wider goal of keeping
privileged code reviewable.

\subsubsection{Block Device Abstraction}

The file system never issues port I/O itself. It operates on a
\texttt{block\_device\_t}, a structure holding a sector size, a sector count,
an opaque context pointer and a table of two function pointers, \texttt{read}
and \texttt{write}, both taking a starting Logical Block Address (LBA) and a
sector count. Two backends implement that interface: an in-memory backend used
during early development, before any disk driver existed, and the ATA/IDE
driver used in the running system. Because the FAT32 code was written against
the abstraction rather than against the driver, the two were developed and
tested independently.

The driver operates the primary and secondary ATA channels in Programmed
Input/Output (PIO) mode. At initialisation it probes up to four
channel and drive combinations, issuing \texttt{IDENTIFY} (\texttt{0xEC}) to
each and accepting the first that answers; the sector count is taken from the
28-bit field of the identify block, falling back to the 48-bit field when that
field reads zero. Reads issue one \texttt{READ SECTORS} (\texttt{0x20}) command
per run of up to 256 sectors, which is the architectural maximum of the sector
count register, and then service one Data Request (DRQ) per sector inside the
run. Batching the command in this way was a deliberate optimisation: the
command itself costs six port writes and a busy-wait for the BSY flag to clear,
and that fixed cost was previously being paid once per sector. Writes are
issued one sector at a time and are followed by a \texttt{CACHE FLUSH}
(\texttt{0xE7}) so that data is not left sitting in the drive's write cache.

The driver validates its own transfer buffers before touching the hardware.
Every page spanned by the buffer is checked against the page tables, and a
transfer whose buffer is not fully mapped is refused rather than attempted.
This is a hard requirement rather than a defensive nicety: the transfer loop is
inside a device command sequence, and a page fault taken part-way through such
a sequence cannot be serviced and resumed cleanly.

\subsubsection{Mounting and Validation}

Mounting reads the boot sector at the start of the volume and rejects it unless
every one of the following holds: the boot signature is \texttt{0xAA55}, the
sector size is 512 bytes, the sectors-per-cluster field is non-zero, the root
directory cluster is at least 2, the FAT size field is non-zero, the file system
type string begins with \texttt{FAT32}, and the version field is zero. From the
accepted boot sector the mount derives the FAT start LBA, the data area start
LBA, the number of data clusters and the cluster size in bytes. The FSInfo
sector is then parsed and accepted only if all three of its signatures
(\texttt{0x41615252}, \texttt{0x61417272}, \texttt{0xAA550000}) are correct, and
the backup boot sector, when the volume declares one, is read and compared
field by field against the primary. A volume that fails any check is not
mounted at all.

The rationale is that every later operation computes disk addresses from these
few fields. Validating them once at mount time means the arithmetic in the
layers above can be written without repeating the same defensive checks, and a
malformed or hostile image is rejected at the one point where rejection is
still cheap.

\subsubsection{Cluster Chains and Allocation}

The data area is addressed in clusters, numbered from 2, and the mapping from a
cluster number to its first sector is

\[
    \mathrm{sector}(c) = \mathrm{data\_start} + (c - 2) \times \mathrm{sectors\_per\_cluster}.
\]

Each FAT entry is 32 bits of which the low 28 are significant, and a value of
\texttt{0x0FFFFFF8} or above marks the end of a chain. Reading an entry
therefore means reading the FAT sector that contains it, and this is where the
first performance problem appeared: a 512-byte sector holds 128 FAT entries, so
walking a chain naively re-reads the same sector 128 times in succession, and
each of those reads is a polled PIO transfer rather than a memory access.
Reading a 4~MB file with 512-byte clusters walks 8192 entries, which is 8192
disk reads.

The mount therefore carries a one-sector FAT cache. One entry is sufficient
because chain walks are sequential by nature: the FAT sector holding the entry
for cluster $N$ is almost always the one holding the entry for cluster $N+1$.
With the cache, the same 8192-entry walk costs 64 sector reads. The cache is
marked invalid before the read that fills it and only marked valid once the
data is known good, so that a failed read cannot leave a torn mixture of two
sectors to be served as if it were a FAT. The write path does not attempt to
patch the cached copy; it invalidates it, on the grounds that FAT writes are
rare (allocating a cluster, extending a file) while the read path they would
complicate is the hot one.

Allocation scans the FAT from cluster 2 for the first entry reading zero,
claims it by writing an end-of-chain marker, and updates the FSInfo free count
and next-free hint. Every FAT write is mirrored into all FAT copies the volume
declares, so the redundant FAT that the format specifies is genuinely
maintained rather than left to drift. Freeing a chain walks it entry by entry,
zeroing each and incrementing the free count, and carries a guard counter
bounded by the total cluster count so that a corrupted volume containing a
cyclic chain terminates with an error instead of looping forever.

\subsubsection{Directories, Paths and Files}

A directory is a cluster chain of 32-byte entries, 16 to a sector. Iteration
walks the chain, skipping entries marked deleted (\texttt{0xE5}) and long file
name entries (attribute \texttt{0x0F}), and stops at the first entry whose name
begins with a zero byte, since that marks the end of the directory rather than
a gap within it. Names are handled in the 8.3 short form: user-supplied names
are converted to upper-case ASCII and split into an eight character name and a
three character extension before comparison, so that lookup does not depend on
the case a program happened to use. On top of iteration sit the four operations
the layers above need: find an entry by name, create one in the first free
slot, update the first cluster and size of an existing one, and delete one.

Path resolution splits an absolute path into components and walks them from the
root cluster, requiring every component except the last to carry the directory
attribute. Opening a file resolves its path, records the parent directory
cluster and short name (needed later to write metadata back), and initialises a
read position of zero.

The read path received the second significant optimisation. Two properties of
real access patterns are exploited. First, reads are usually sequential, so
each open file carries a chain hint: the cluster last reached and its index
within the chain, allowing the next read to resume the walk from there instead
of restarting at the first cluster. The hint cannot silently go stale, because
the only operation that rewrites a chain is truncation, which resets it, while
appending only adds indices beyond any the hint can name. Second, files on a
freshly written image are laid out contiguously, so before issuing a transfer
the reader looks ahead through the FAT to find how many clusters continue
physically next door and reads that whole run in one command. For the 4~MB WAD
file used by the ported game, which occupies 8192 consecutive clusters, this
collapses 8192 single-sector commands into 128 multi-sector ones. A fragmented
file is not a special case; it simply yields shorter runs.

Sectors are never read directly into the caller's buffer. They are read into a
kernel bounce buffer, allocated at eight sectors and falling back to a single
512-byte stack buffer when the heap allocation fails, and the requested bytes
are copied out from there. This is required by the isolation model rather than
by convenience: the destination may be a userland pointer whose pages are
demand-paged and not yet present, and the driver refuses unmapped destinations
because it cannot service a fault from inside a transfer.

Writing is a read-modify-write at sector granularity: the affected sector is
read, the new bytes are copied in, and the sector is written back, with clusters
allocated and linked on demand as the file grows past its current chain.
Truncation frees the tail of the chain, resets the read hint and rewrites the
directory entry; deletion frees the whole chain and clears the entry.

\subsubsection{Exposure to User Programs}

User programs reach the file system only through system calls. Each process
carries its own descriptor table of sixteen slots, in which the first three are
reserved for standard input, output and error, and the fourth is pinned by the
program loader to the executable the process is running. That pinned descriptor
behaves like any other open file for reading and seeking, but \texttt{close}
refuses it, because every page of program text and data that has not yet been
touched is demand-paged through it for the lifetime of the process; closing it
from user space would remove the backing store the calling code is executing
from.

Because descriptor tables are per-process, one process cannot name another
process's open files, and the descriptor numbers a program sees are indices
into its own table rather than into any global structure. Every pointer that
arrives from user space is validated against the calling process's own address
space before it is dereferenced: buffers are checked as a range, and path
strings are checked as a bounded string with a 256 byte ceiling. A pointer that
fails validation causes the call to return an error, so that a bad or malicious
pointer from ring 3 produces a failed system call rather than a kernel fault.

\subsubsection{Verification}

Correctness was established by three complementary means rather than by
inspection alone.

The first is a self-test suite that runs inside the kernel at every boot and
reports each result over the serial line. It covers resolving and listing the
root directory, resolving and reading a file in the root and a file in a
subdirectory, and comparing the returned bytes against known content; it then
performs a full write cycle on a scratch file, creating it, writing, appending,
reading the result back and comparing it, truncating it, deleting it and
confirming that resolution afterwards fails. Mount-level invariants such as the
root cluster number and the FSInfo signatures are asserted as well. Because
these run on every boot, a regression in the file system is visible immediately
rather than at the point some later feature happens to depend on it.

The second is the independence of the test data already noted: the images are
formatted and populated by host tools, so passing the suite means agreeing with
an implementation written by someone else.

The third is end-to-end use by real programs. A user space \texttt{cat} utility
reads, seeks within and checksums the 4~MB WAD file through the C library's
buffered I/O, exercising the whole stack from newlib down to the PIO driver;
and the ported game reads that same file at startup through ordinary system
calls, which makes it an integration test of the read path, the chain hint, the
contiguous run detection and the bounce buffer under sustained load.
